\section{Bất đẳng thức và tính chất}

\begin{lesson}
    \subsection{Kiến thức trọng tâm}

    \subsubsection{Bất đẳng thức}

    \textbf{Nhắc lại thứ tự trên tập số thực}

    \begin{keyconcept}
        Trên tập số thực, với hai số $a$ và $b$ có ba trường hợp sau:
        \begin{enumerate}[label=\alph*)]
            \item Số $a$ bằng số $b$, kí hiệu $a=b$;
            \item Số $a$ lớn hơn số $b$, kí hiệu $a>b$;
            \item Số $a$ nhỏ hơn số $b$, kí hiệu $a<b$.
        \end{enumerate}
    \end{keyconcept}

    \textbf{Câu hỏi}

    Thay ? trong các biểu thức sau bằng dấu thích hợp $(=,>,<)$.
    \begin{enumerate}[label=\alph*)]
        \item $-34{,}2\ \answerbox\ -27$;
        \item $\dfrac{6}{-8}\ \answerbox\ -\dfrac{3}{4}$;
        \item $2\,024\ \answerbox\ 1\,954$.
    \end{enumerate}

    \begin{keyconcept}
        Khi biểu diễn số thực trên trục số, điểm biểu diễn số bé hơn nằm trước điểm biểu diễn số lớn hơn. Chẳng hạn, $-2{,}5<-1<1<1{,}5$.

        Số $a$ lớn hơn hoặc bằng số $b$, tức là $a>b$ hoặc $a=b$, kí hiệu là $a\ge b$.

        Số $a$ nhỏ hơn hoặc bằng số $b$, tức là $a<b$ hoặc $a=b$, kí hiệu là $a\le b$.
    \end{keyconcept}

    \textbf{Khái niệm bất đẳng thức}

    \begin{keyconcept}
        Ta gọi hệ thức dạng $a>b$ (hay $a<b$, $a\ge b$, $a\le b$) là bất đẳng thức và gọi $a$ là vế trái, $b$ là vế phải của bất đẳng thức.
    \end{keyconcept}

    \textbf{Chú ý:} Hai bất đẳng thức $1<2$ và $-3<-2$ (hay $6>3$ và $8>5$) được gọi là hai bất đẳng thức cùng chiều. Hai bất đẳng thức $1<2$ và $-2>-3$ (hay $6>3$ và $5<8$) được gọi là hai bất đẳng thức ngược chiều.

    \begin{keyconcept}
        Bất đẳng thức có tính chất quan trọng sau:

        Nếu $a<b$ và $b<c$ thì $a<c$ (tính chất bắc cầu của bất đẳng thức).
    \end{keyconcept}

    \textbf{Chú ý:} Tương tự, các thứ tự lớn hơn $(>)$, lớn hơn hoặc bằng $(\ge)$, nhỏ hơn hoặc bằng $(\le)$ cũng có tính chất bắc cầu.

    \textbf{Luyện tập 2}

    Chứng minh rằng:
    \begin{enumerate}[label=\alph*)]
        \item $\dfrac{2\,024}{1\,000}>1{,}9$;

        \answerline
        \item $-\dfrac{2\,022}{2\,023}>-1{,}1$.

        \answerline
    \end{enumerate}

    \subsubsection{Liên hệ giữa thứ tự và phép cộng}

    \textbf{Khám phá}

    Xét bất đẳng thức $-1<2$. \hfill (1)
    \begin{enumerate}[label=\alph*)]
        \item Cộng $2$ vào hai vế của bất đẳng thức (1) rồi so sánh kết quả thì ta được bất đẳng thức nào?

        \answerline
        \item Cộng $-2$ vào hai vế của bất đẳng thức (1) rồi so sánh kết quả nhận được thì ta được bất đẳng thức nào?

        \answerline
        \item Cộng vào hai vế của bất đẳng thức (1) với cùng một số $c$ thì ta sẽ được bất đẳng thức nào?

        \answerline
    \end{enumerate}

    \begin{keyconcept}
        Khi cộng cùng một số vào hai vế của một bất đẳng thức ta được bất đẳng thức mới cùng chiều với bất đẳng thức đã cho.

        Với ba số $a$, $b$, $c$ ta có:
        \begin{itemize}
            \item Nếu $a<b$ thì $a+c<b+c$;
            \item Nếu $a\le b$ thì $a+c\le b+c$;
            \item Nếu $a>b$ thì $a+c>b+c$;
            \item Nếu $a\ge b$ thì $a+c\ge b+c$.
        \end{itemize}
    \end{keyconcept}

    \textbf{Luyện tập 3}

    Không thực hiện phép tính, hãy so sánh:
    \begin{enumerate}[label=\alph*)]
        \item $19+2\,023$ và $-31+2\,023$;

        \answerline
        \item $\sqrt{2}+2$ và $4$.

        \answerline
    \end{enumerate}

    \subsubsection{Liên hệ giữa thứ tự và phép nhân}

    \textbf{Khám phá}

    Cho bất đẳng thức $-2<5$.
    \begin{enumerate}[label=\alph*)]
        \item Nhân hai vế của bất đẳng thức với $7$ rồi so sánh kết quả thì ta được bất đẳng thức nào?

        \answerline
        \item Nhân hai vế của bất đẳng thức với $-7$ rồi so sánh kết quả thì ta được bất đẳng thức nào?

        \answerline
    \end{enumerate}

    \begin{keyconcept}
        Khi nhân cả hai vế của một bất đẳng thức với cùng một số dương ta được bất đẳng thức mới cùng chiều với bất đẳng thức đã cho.

        Với ba số $a$, $b$, $c$ và $c>0$ ta có:
        \begin{itemize}
            \item Nếu $a<b$ thì $ac<bc$;
            \item Nếu $a\le b$ thì $ac\le bc$;
            \item Nếu $a>b$ thì $ac>bc$;
            \item Nếu $a\ge b$ thì $ac\ge bc$.
        \end{itemize}
    \end{keyconcept}

    \begin{keyconcept}
        Khi nhân cả hai vế của một bất đẳng thức với cùng một số âm ta được bất đẳng thức mới ngược chiều với bất đẳng thức đã cho.

        Với ba số $a$, $b$, $c$ và $c<0$ ta có:
        \begin{itemize}
            \item Nếu $a<b$ thì $ac>bc$;
            \item Nếu $a\le b$ thì $ac\ge bc$;
            \item Nếu $a>b$ thì $ac<bc$;
            \item Nếu $a\ge b$ thì $ac\le bc$.
        \end{itemize}
    \end{keyconcept}

    \textbf{Luyện tập 4}

    Thay ? trong các biểu thức sau bởi dấu thích hợp $(<,>)$ để được khẳng định đúng.
    \begin{enumerate}[label=\alph*)]
        \item $13\cdot(-10{,}5)\ \answerbox\ 13\cdot 11{,}2$;
        \item $(-13)\cdot(-10{,}5)\ \answerbox\ (-13)\cdot 11{,}2$.
    \end{enumerate}

    \textbf{Vận dụng 2}

    Một nhà tài trợ dự kiến tổ chức một buổi đi dã ngoại tập thể nhằm giúp các bạn học sinh vùng cao trải nghiệm thực tế tại một trang trại trong 1 ngày (từ 14h00 ngày hôm trước đến 12h00 ngày hôm sau). Cho biết số tiền tài trợ dự kiến là 30 triệu đồng và giá thuê các dịch vụ và phòng nghỉ là 17 triệu đồng 1 ngày, giá mỗi suất ăn trưa, ăn tối là 60 000 đồng và mỗi suất ăn sáng là 30 000 đồng. Hỏi có thể tổ chức cho nhiều nhất bao nhiêu bạn tham gia được?

    \answerline
    \answerline
    \answerline
\end{lesson}

\begin{exercises}
    \subsection{Bài tập}

    \begin{questions}[resume=SGK]
        \item Dùng kí hiệu để viết bất đẳng thức tương ứng trong mỗi trường hợp sau:
        \begin{questions}
            \item $x$ nhỏ hơn hoặc bằng $-2$;
            \item $m$ là số âm;
            \item $y$ là số dương;
            \item $p$ lớn hơn hoặc bằng $2\,024$.
        \end{questions}
        \answerline

        \item Viết một bất đẳng thức phù hợp trong mỗi trường hợp sau:
        \begin{questions}
            \item Bạn phải ít nhất 18 tuổi mới được phép lái ô tô;
            \item Xe buýt chở được tối đa 45 người;
            \item Mức lương tối thiểu cho một giờ làm việc của người lao động là 20 000 đồng.
        \end{questions}
        \answerline

        \item Không thực hiện phép tính, hãy chứng minh:
        \begin{questions}
            \item $2\cdot(-7)+2\,023<2\cdot(-1)+2\,023$;

            \answerline
            \item $(-3)\cdot(-8)+1\,975>(-3)\cdot(-7)+1\,975$.

            \answerline
        \end{questions}

        \item Cho $a<b$, hãy so sánh:
        \begin{questions}
            \item $5a+7$ và $5b+7$;

            \answerline
            \item $-3a-9$ và $-3b-9$.

            \answerline
        \end{questions}

        \item So sánh hai số $a$ và $b$, nếu:
        \begin{questions}
            \item $a+1\,954<b+1\,954$;

            \answerline
            \item $-2a>-2b$.

            \answerline
        \end{questions}

        \item Chứng minh rằng:
        \begin{questions}
            \item $-\dfrac{2\,023}{2\,024}>-\dfrac{2\,024}{2\,023}$;

            \answerline
            \item $\dfrac{34}{11}>\dfrac{26}{9}$.

            \answerline
        \end{questions}
    \end{questions}
\end{exercises}

\begin{practice}
    \subsection{Luyện tập}

    \begin{questions}[resume=SBT]
        % TODO(HINH-NGUON): bo sung asset/crop dung tu SBT PDF trang 25 cho Bai 2.7; khong redraw xap xi bang TikZ.
        \item Khi đi trên tuyến cao tốc Thành phố Hồ Chí Minh -- Trung Lương, chúng ta thấy biển báo giao thông báo hiệu giới hạn tốc độ mà ô tô được phép đi trong điều kiện bình thường. Hãy viết các bất đẳng thức để mô tả tốc độ cho phép của ô tô
        \begin{questions}
            \item Ở làn ngoài cùng bên trái và ở làn giữa;
            \item Ở làn ngoài cùng bên phải.
        \end{questions}
        \answerline

        \item Viết bất đẳng thức để mô tả tình huống sau:
        \begin{questions}
            \item Bạn phải ít nhất 18 tuổi mới được đi bầu cử đại biểu Quốc hội.
            \item Một thang máy chở được tối đa 700 kg.
            \item Bạn phải mua hàng có tổng giá trị ít nhất 1 triệu đồng mới được giảm giá.
            \item Bạn phải ném vào rổ ít nhất 5 quả bóng mới vào được đội tuyển bóng rổ.
        \end{questions}
        \answerline

        \item So sánh
        \begin{questions}
            \item $-\dfrac{2\,019}{1\,010}$ và $-\dfrac{201}{100}$;

            \answerline
            \item $\dfrac{2\,024^2-1}{2\,024}$ và $\dfrac{2\,025^2+1}{2\,025}$.

            \answerline
        \end{questions}

        \item Cho $a>b$, hãy so sánh
        \begin{questions}
            \item $20a+5b$ và $20b+5a$;

            \answerline
            \item $-3(a+b)-1$ và $-6b-1$.

            \answerline
        \end{questions}

        \item Cho $a>b>0$, chứng minh rằng
        \begin{questions}
            \item $a^2>ab$ và $ab>b^2$;

            \answerline
            \item $a^2>b^2$ và $a^3>b^3$.

            \answerline
        \end{questions}

        \textbf{Chú ý:} Tính chất ``Với $a>b>0$ thì $a^2>b^2$ và $a^3>b^3$'' thường hay dùng trong nhiều bài toán chứng minh bất đẳng thức.

        \item Chứng minh rằng với mọi số $a$, $b$ ta có $\dfrac{a^2+b^2}{2}\ge ab$.

        \answerline
        \answerline

        \item Chứng minh rằng: Trong ba số tự nhiên liên tiếp thì bình phương số đứng giữa lớn hơn tích hai số còn lại.

        \answerline
        \answerline
    \end{questions}
\end{practice}
